% !TEX TS-program = xelatex
% !TEX encoding = UTF-8 Unicode
% !Mode:: "TeX:UTF-8"

\documentclass{resume}
\usepackage{zh_CN-Adobefonts_external} % Simplified Chinese Support using external fonts (./fonts/zh_CN-Adobe/)
%\usepackage{zh_CN-Adobefonts_internal} % Simplified Chinese Support using system fonts
\usepackage{linespacing_fix} % disable extra space before next section
\usepackage{cite}

\begin{document}
\pagenumbering{gobble} % suppress displaying page number

\name{张志锐(ZHIRUI ZHANG)}
% \name{越空}

\basicInfo{
  \email{zrustc11@gmail.com} \textperiodcentered\ 
  \phone{(+86) 188-106-36928} \textperiodcentered\ 
  \text{\faFilesO\ https://zrustc.github.io/}
  }

\section{\faGraduationCap\  教育背景}
\datedsubsection{\textbf{中国科学技术大学}}{2014 -- 2019}
\textit{中国科学技术大学-微软亚洲研究院联合培养博士}\ \ 计算机科学与技术 \\
\textit{导师: 沈向洋、陈恩红 \ \ 副导师：李沐} \\
\textit{研究方向：深度学习、自然语言处理、机器翻译、自然语言生成} \\
\text{Google Scholar Citations: $2400^{+}$,\ h-index: 20+}
\datedsubsection{\textbf{中国科学技术大学}}{2010 -- 2014}
\textit{本科}\ \ 计算机科学与技术


\textit{}

\section{\faUsers\ 主要经历}

\datedsubsection{\textbf{华为}，深圳}{2024.3 -- }
\role{终端-小艺业务部，技术专家A (20级)，从零构建团队8人（1个天少），1次A绩效}
\text{主要负责大模型通用后训练、慢思考推理模型以及支撑模型场景落地}
\begin{itemize}
\item 多语言大模型构建和落地 (2024.3-2024.8)：(1) 负责小艺多语言大模型构建，牵引诺亚实现盘古多语言模型能力提升，包括高效多语言词表扩充（压缩率分析），多语言篇章数据挖掘和过滤，多语言训练策略数据配比验证，多语言指令对齐数据构建等；(2) 基于多语言大模型进行场景微调并落地小艺翻译，探索1.5B-LLM端侧大模型结合语音输入数据增强策略，实现中英双向端到端语音翻译性能平均提升7个COMET值，体感显著优于三星手机翻译30\%+；7B-LLM中英互译能力在内部测试集上显著优于云侧最优NMT模型，实现公开翻译测试集上持平GPT-4性能；(3) 协助整个小艺翻译团队梳理整个技术框架和方案；
\item 大模型通用后训练 (2024.9-now)：(1) 负责小艺通用后训练阶段，梳理盘古通用后训练流程和问题，牵引诺亚实现SFT3.0构建，包括指令多样性体系、高效自评测体系，两阶段通用SFT训练策略，支撑通用数据和模型高效迭代，迭代速度从周级别降低到天级别；(2) Qwen2.5-7B基座模型+这套训练策略，在主流开源测试集（知识、推理、数学、代码、指令遵从、工具调用）均超过Qwen2.5-7B-Instruct，其中AlpacaEval 2.0和Arena-hard提升10\%+, 内部测试集各个场景提升5\%--10\%；相同策略在盘古模型38B/135B上实现10\%+分位值提升；(3) 相应策略应用于多模态大模型，近期基于Qwen3-8B构建多模态大模型+后训练高质量指令数据退火+SFT3.0训练策略，实现Qwen3-8B-VL的通用能力持平/超过Qwen2.5-32B-VL（数学和学科类持平、指令遵从+5\%、ArenaHard持平等）；(4) 支撑430小艺Live态对话亮点特性发布和上线；RL+GenRM训练策略提升小艺Live态对话满意度；
\item 慢思考推理模型 (2024.9-now)：(1) 负责通用慢思考推理模型构建以及前沿技术探索，完成多项技术探索和验证，包括BoN+PRM能力边界、快慢思考融合+可控思考长度策略（早于Qwen3发布）、解题和判别能力的双生特性、基于置信度的RL优化算法、环境交互RL等，\textbf{推演出推理模型的第一性原理}、完成OpenAI的O系列low/medium/high/pro思考模式的技术方案解密；(2) 小艺深度解题智能体技术leader，落地面向教育场景的快慢思考融合推理模型、AI交互讲解亮点特性打造等，在2025电信打榜实现AI学习能力第一；(3) 探索基于LLM的Live2D数字人驱动策略方案，实现第一版本基于大模型驱动的离线可用素材库生成；
\end{itemize}

\datedsubsection{\textbf{阶跃星辰}，北京}{2023.11 -- 2024.3}
\text{算法专家，主要负责语言大模型预训练工作以及前沿技术探索}
\begin{itemize}
\item Step2万亿大模型预训练项目核心技术骨干，协助验证数据和MoE结构等相关算法策略；探索和落地FP8-LM训练技术（包括预训练和后训练阶段）；
\item 探索长序列策略，验证各类主流LongContext策略，支持128K上下文落地；
\end{itemize}

\datedsubsection{\textbf{腾讯-AI Lab}，深圳}{2021.9 -- 2023.11}
\role{技术级别: T11, 自然语言算法平台组，1次5星, 2次SEVP卓越个人}
\text{主要负责腾讯翻译训练平台搭建、交互翻译模型建设、个性化机器翻译以及先进技术研究}
\begin{itemize}
  \item 负责腾讯翻译分布式训练平台搭建和优化，自动化训练流程搭建（包括高质量翻译数据过滤和领域识别），交互翻译多合一模型(具备自动翻译、约束解码和翻译记忆融合功能)构建与优化，个性化机器翻译策略研发与落地，参与翻译商业化项目；协助管理自然语言算法平台组翻译方向人员；
  \item 开展自然语言处理相关研究，提出联邦机器翻译,交互机器翻译开源平台, Simple and Scalable Nearest Neighbor Machine Translation (SK-MT)等新方法，其中\textbf{SK-MT成为腾讯翻译的个性化机器翻译核心策略并完成上线}；另外探索模型鲁棒性相关方法，包括理解词级别文本攻击和通过Nearest Neighbor Calibration提升LLM的In-Context Learning能力; 
  \item 开展多模态相关机器翻译探索，包括Speech-to-Speech Translation和Automatic Song Translation等方向，取得IWSLT2023比赛中\textbf{端到端语音到语音翻译子任务第一名}；
  \item 参与混元大模型项目，负责多语言预训练数据清洗，LLM的翻译能力和RLHF相关调优；
\end{itemize}

\datedsubsection{\textbf{阿里巴巴-达摩院}，杭州}{2019.7 -- 2021.8}
\role{算法专家，语言技术实验室，阿里星，1次375绩效}
% {主管：陈博兴、谢军}
\text{主要负责阿里翻译基础通用模型建设、语音翻译、对外商业化以及先进技术研究}
\begin{itemize}
  \item 负责阿里云通用翻译服务以及商业化项目，技术上完成阿里通用翻译自动化链路构建、通用模型质量优化（中英互译内部评测达到Top1）、通用语向扩展（15->214）、线上模型解码加速（2-3x)，以及商业化领域模型定制和交付;
  \item 开展自然语言处理相关研究，提出高效的多语言机器翻译训练策略，融合预训练语言模型的机器翻译，Iterative Domain-Repaired Back-Translation，端到端语音翻译模型，以及初步探索个性化机器翻译策略。这些方法已用于大幅度提升线上多语言翻译模型性能和领域模型定制效果。另外也探索可变长的文本对抗攻击和任务型对话系统作为自然语言生成等新方法;
  \item 受邀在DataFunTalk 2020年终大会上分享“\textbf{阿里多语言翻译模型的前沿探索及技术实践}”的主题报告;
\end{itemize}

\datedsubsection{\textbf{微软亚洲研究院}，北京}{2019年2月 -- 2019年6月}
\role{研究实习生，自然语言计算组，(联合培养博士生项目)}{指导老师: 刘树杰}
\textit{自然语言处理相关技术研究}

\datedsubsection{\textbf{微软雷德蒙德研究院}，西雅图，美国}{2018年7月 -- 2019年1月}
\role{Research Intern，Deep Learning Group,}{指导老师: Xiujun Li, Jianfeng Gao}
\textit{对话系统相关技术研究}
\begin{itemize}
  \item 在资源受限的任务型对话场景下，提出Budgeted Policy Learning方法来显著提升对话系统的成功率；
\end{itemize}

\datedsubsection{\textbf{微软亚洲研究院}，北京}{2015年7月 -- 2018年7月}
\role{研究实习生，自然语言计算组，(联合培养博士生项目)}{指导老师: 李沐}
\textit{机器翻译、深度学习、自然语言生成等相关技术研究}
\begin{itemize}
  \item 开展神经机器翻译相关研究，提出Iterative Back-Translation, Target-bidirectional Agreement, Coarse-to-Fine Learning, Bidirectional GANs等训练算法，针对无监督机器翻译设计SMT as Posterior Regularization方法，另外将NMT相关技术应用于Dependency Parsing和Text Style Transfer任务中;
  \item 支持微软内部项目开发(SmartFlow Toolkit、Writing Intelligence和Babel Project):
  \item[-] \emph{SmartFlow Toolkit (2016)}:  开发C\#版本的深度学习工具，实现OP计算、计算图调度，内部管理等功能，并服务于组内各类模型开发（包括R-NET模型、对话/QA相关的NLG模型等），部分模型已应用于微软小冰对话生成和必应产品中;
  \item[-] \emph{Writing Intelligence Project (2017)}: 该项目想借助深度学习技术去使写作更便捷，其中定义了句子补全、基于关键词的句子生成、语法检查、下一句话推荐这四个功能。主要负责句子补全功能的模型开发以及实现;
  \item[-] \emph{Babel Project (2017-2018)}: 微软内部多个小组一起合作，探索在大规模语料下当前神经机器翻译的极限。负责将Iterative Back-Translation和Target-bidirectional Agreement两个技术应用于该项目，显著提升模型的性能，并达到人类“可比”水平;
\end{itemize}

\datedsubsection{\textbf{微软亚洲研究院}，北京}{2013年7月 -- 2014年6月}
\role{研究实习生，自然语言计算组}{指导老师: 李沐}
\textit{自然语言处理相关技术研究}
\begin{itemize}
  \item 熟悉和复现各类基础NLP模型，将word2vec技术用于篇章关系分析中；
\end{itemize}

\section{\faFilesO\ 主要论文}

\begin{enumerate}

\item \normalsize \textbf{Document-Level Machine Translation with Large Language Models}\\ 
\small - Longyue Wang*, Chenyang Lyu*, Tianbo Ji*, \textbf{Zhirui Zhang*}, Dian Yu, Shuming Shi, Zhaopeng Tu. In \textit{EMNLP}'2023.

\item \normalsize \textbf{Nearest Neighbor Machine Translation is Meta-Optimizer on Output Projection Layer}\\ 
\small - Ruize Gao,  \textbf{Zhirui Zhang}, Yichao Du, Lemao Liu, Rui Wang. In \textit{EMNLP}'2023.

\item \normalsize \textbf{IMTLab: An Open-Source Platform for Building, Evaluating, and Diagnosing Interactive Machine Translation Systems}\\ 
\small - Xu Huang, \textbf{Zhirui Zhang}, Ruize Gao, Yichao Du, Lemao Liu, Guoping Huang, Shuming Shi, Jiajun Chen, Shujian Huang. In \textit{EMNLP}'2023.

\item \normalsize \textbf{Rethinking Word-Level Auto-Completion in Computer-Aided Translation}\\ 
\small - Xingyu Chen, Lemao Liu, Guoping Huang, \textbf{Zhirui Zhang}, Mingming Yang, Shuming Shi, Rui Wang. In \textit{EMNLP}'2023.

\item \normalsize \textbf{Fairness-guided Few-shot Prompting for Large Language Models}\\ 
\small - Huan Ma, Changqing Zhang, Yatao Bian, Lemao Liu, \textbf{Zhirui Zhang}, Peilin Zhao, Shu Zhang, Huazhu Fu, Qinghua Hu, Bingzhe Wu. In \textit{NeurIPS}'2023.

\item \normalsize \textbf{The \textsc{MineTrans} Systems for IWSLT 2023 Offline Speech Translation and Speech-to-Speech Translation Tasks}\\ 
\small - Yichao Du, Zhengsheng Guo, Jinchuan Tian, \textbf{Zhirui Zhang}, Xing Wang, Jianwei Yu, Zhaopeng Tu, Tong Xu and Enhong Chen. In \textit{IWSLT}'2023.

\item \normalsize \textbf{Rethinking Translation Memory Augmented Neural Machine Translation}\\
\small -  Hongkun Hao, Guoping Huang, Lemao Liu, \textbf{Zhirui Zhang}, Shuming Shi and Rui Wang. In \textit{ACL}'2023.

\item \normalsize \textbf{E-NER: Evidential Deep Learning for Trustworthy Named Entity Recognition}\\
\small - Zhen Zhang, Mengting Hu, Shiwan Zhao, Minlie Huang, Haotian Wang, Lemao Liu, \textbf{Zhirui Zhang}, Zhe Liu and Bingzhe Wu.  In \textit{ACL}'2023.

\item \normalsize \textbf{Federated Nearest Neighbor Machine Translation}\\
\small - Yichao Du, \textbf{Zhirui Zhang}, Bingzhe Wu, Lemao Liu, Tong Xu and Enhong Chen. In \textit{ICLR}'2023.

\item \normalsize \textbf{Simple and Scalable Nearest Neighbor Machine Translation}\\
\small - Yuhan Dai*, \textbf{Zhirui Zhang*}, Qiuzhi Liu, Qu Cui, Weihua Li, Yichao Du and Tong Xu. In \textit{ICLR}'2023.

\item \normalsize \textbf{Context-Adaptive Document-Level Neural Machine Translation}\\
\small - Linlin Zhang, \textbf{Zhirui Zhang}, Boxing Chen, Weihua Luo and Lou Si. In \textit{ICASSP}'2022.

\item \normalsize \textbf{Non-Parametric Domain Adaptation for End-to-End Speech Translation}\\
\small - Yichao Du*, Weizhi Wang*, \textbf{Zhirui Zhang}, Boxing Chen, Jun Xie, Tong Xu and Enhong Chen. In \textit{EMNLP}'2022.

\item \normalsize \textbf{Task-Oriented Dialogue System as Natural Language Generation}\\
\small - Weizhi Wang, \textbf{Zhirui Zhang}, Junliang Guo, Yinpei Dai, Boxing Chen and Weihua Luo. In \textit{SIGIR}'2022.

\item \normalsize \textbf{Automatic Song Translation for Tonal Languages}\\
\small - Fenfei Guo, Chen Zhang, \textbf{Zhirui Zhang}, Qixin He, Kejun Zhang, Jun Xie and Jordan Lee Boyd-Graber. In \textit{ACL}'2022.


\item \normalsize \textbf{Regularizing End-to-End Speech Translation with Triangular Decomposition Agreement}\\
\small - Yichao Du, \textbf{Zhirui Zhang}, Weizhi Wang, Boxing Chen, Jun Xie and Tong Xu. In \textit{AAAI}'2022.

\item \normalsize \textbf{Non-Parametric Online Learning from Human Feedback for Neural Machine Translation}\\
\small - Dongqi Wang, Hao-Ran Wei, \textbf{Zhirui Zhang}, Shujian Huang, Jun Xie, Weihua Luo and Jiajun Chen. In \textit{AAAI}'2022.

\item \normalsize \textbf{Non-Parametric Unsupervised Domain Adaptation for Neural Machine Translation}\\
\small - Xin Zheng, \textbf{Zhirui Zhang}, Shujian Huang, Boxing Chen, Jun Xie, Weihua Luo and Jiajun Chen. In \textit{EMNLP}'2021.

\item \normalsize \textbf{Rethinking Zero-shot Neural Machine Translation: From a Perspective of Latent Variables}\\
\small - Weizhi Wang, \textbf{Zhirui Zhang}, Yichao Du, Boxing Chen, Jun Xie and Weihua Luo. In \textit{EMNLP}'2021.

\item \normalsize \textbf{Sentence-State LSTMs for Sequence-to-Sequence Learning}\\
\small - Xuefeng Bai, Yafu Li, \textbf{Zhirui Zhang}, Mingzhou Xu, Boxing Chen, Weihua Luo, Derek Wong and Yue Zhang. In \textit{NLPCC}'2021.


\item \normalsize \textbf{Adaptive Nearest Neighbor Machine Translation}\\
\small - Xin Zheng, \textbf{Zhirui Zhang}, Junliang Guo, Shujian Huang, Boxing Chen, Weihua Luo and Jiajun Chen. In \textit{ACL-IJCNLP}'2021.

\item \normalsize \textbf{Adaptive Adapter: an Efficient
Way to Incorporate BERT into Neural Machine Translation}\\
\small - Junliang Guo, \textbf{Zhirui Zhang}, Linli Xu, Boxing Chen and Enhong Chen. In \textit{TASLP}'2021.

\item \normalsize \textbf{Incorporating BERT into Parallel Sequence Decoding with Adapters}\\
\small - Junliang Guo, \textbf{Zhirui Zhang}, Linli Xu, Hao-Ran Wei, Boxing Chen and Enhong Chen. In \textit{NeurIPS}'2020

\item \normalsize \textbf{Iterative Domain-Repaired Back-Translation} \\
\small - Hao-Ran Wei, \textbf{Zhirui Zhang}, Boxing Chen and Weihua Luo. In \textit{EMNLP}'2020.

\item \normalsize \textbf{Cross-lingual Pre-training Based Transfer for Zero-shot Neural Machine Translation} \\
\small - Baijun Ji, \textbf{Zhirui Zhang}, Xiangyu Duan, Min Zhang, Boxing Chen and Weihua Luo. In \textit{AAAI}'2020.

\item \normalsize \textbf{Budgeted Policy Learning for Task-Oriented Dialogue Systems} \\
\small - \textbf{Zhirui Zhang}, Xiujun Li, Jianfeng Gao and Enhong Chen. In \textit{ACL}'2019.


\item \normalsize \textbf{Regularizing Neural Machine Translation by Target-bidirectional Agreement} \\
\small - \textbf{Zhirui Zhang}, Shuangzhi Wu, Shujie Liu, Mu Li, Ming Zhou and Tong Xu. In \textit{AAAI}'2019.

\item \normalsize \textbf{Unsupervised Neural Machine Translation with SMT as Posterior Regularization} \\
\small - Shuo Ren*, \textbf{Zhirui Zhang*}, Shujie Liu, Ming Zhou and Shuai Ma. In \textit{AAAI}'2019.

\item \normalsize \textbf{Dependency-to-Dependency Neural Machine Translation} \\
\small - Shuangzhi Wu, Dongdong Zhang, \textbf{Zhirui Zhang}, Nan Yang, Mu Li and Ming Zhou. In \textit{TASLP}'2018.

\item \normalsize \textbf{Bidirectional Generative Adversarial Networks for Neural Machine Translation} \\
\small - \textbf{Zhirui Zhang}, Shujie Liu, Mu Li, Ming Zhou and Enhong Chen. In \textit{CoNLL}'2018.

\item \normalsize \textbf{Coarse-To-Fine Learning for Neural Machine Translation} \\
\small - \textbf{Zhirui Zhang}, Shujie Liu, Mu Li, Ming Zhou and Enhong Chen. In \textit{NLPCC}'2018.

\item \normalsize \textbf{Joint Training for Neural Machine Translation Models with Monolingual Data} \\
\small - \textbf{Zhirui Zhang}, Shujie Liu, Mu Li, Ming Zhou and Enhong Chen. In \textit{AAAI}'2018.

\item \normalsize \textbf{Generative Bridging Network in Neural Sequence Prediction} \\
\small - Wenhu Chen, Guanlin Li, Shuo Ren, Shujie Liu, \textbf{Zhirui Zhang}, Mu Li and Ming Zhou. In \textit{NAACL}'2018.

\item \normalsize \textbf{Learning to Collaborate for Question Answering and Asking} \\
\small - Duyu Tang, Nan Duan, Zhao Yan, \textbf{Zhirui Zhang}, Yibo Sun, Shujie Liu, Yuanhua Lv and Ming Zhou. In \textit{NAACL}'2018.

\item \normalsize \textbf{Stack-based Multi-layer Attention for Transition-based Dependency Parsing} \\
\small - \textbf{Zhirui Zhang}, Shujie Liu, Mu Li, Ming Zhou and Enhong Chen. In \textit{EMNLP}'2017.

\end{enumerate}

\section{\faFilesO\ 技术报告和Preprints}}

\begin{enumerate}

\item \normalsize \textbf{Less is More: Understanding Word-level Textual Adversarial Attack via n-gram Frequency Descend}\\
\small - Ning Lu, Shengcai Liu, \textbf{Zhirui Zhang}, Qi Wang, Haifeng Liu and Ke Tang. \textit{Pre-print}, https://arxiv.org/abs/2302.02568

\item \normalsize \textbf{Improving Few-Shot Performance of Language Models via Nearest Neighbor Calibration}\\
\small - Feng Nie, Meixi Chen, \textbf{Zhirui Zhang} and Xu Cheng. \textit{Pre-print}, https://arxiv.org/abs/2212.02216

\item \normalsize \textbf{Towards Variable-Length Textual Adversarial Attacks}\\
\small - Junliang Guo, \textbf{Zhirui Zhang}, Linlin Zhang, Linli Xu, Boxing Chen, Enhong Chen and Weihua Luo. \textit{Pre-print}, https://arxiv.org/abs/2104.08139 

\item \normalsize \textbf{Style Transfer as Unsupervised Machine Translation} \\
\small - \textbf{Zhirui Zhang*}, Shuo Ren*, Shujie Liu, Jianyong Wang, Peng Chen, Mu Li, Ming Zhou and Enhong Chen. \textit{Pre-print}, https://arxiv.org/abs/1808.07894

\item \normalsize \textcolor{red}{\textbf{Achieving Human Parity on Automatic Chinese to English News Translation}} \\
\small - Hany Hassan, Anthony Aue, Chang Chen, Vishal Chowdhary, Jonathan Clark, Christian Federmann, Xuedong Huang, Marcin Junczys-Dowmunt, William Lewis, Mu Li, Shujie Liu, Tie-Yan Liu, Renqian Luo, Arul Menezes, Tao Qin, Frank Seide, Xu Tan, Fei Tian, Lijun Wu, Shuangzhi Wu, Yingce Xia, Dongdong Zhang, \textbf{Zhirui Zhang} and Ming Zhou. \textit{Pre-print}, https://arxiv.org/abs/1803.05567 

\end{enumerate}

\section{\faCogs\ 主要技能及学术服务}
% increase linespacing [parsep=0.5ex]
\begin{itemize}
  \item 专业技能: 机器翻译/自然语言生成/对话系统/预训练语言模型/对抗攻击/依存句法分析等
  \item 编程技能: Python/C++/C\#/Java/CUDA/Cudnn/Tensorflow/Pytorch
  \item 学术服务: 担任多个NLP/ML会议与期刊长期审稿人，包括ACL、EMNLP、NAACL、AAAI、IJCAI、NeurIPS、ICML、ICLR，COLR等，担任过ACL2024领域主席;
\end{itemize}

\section{\faHeartO\ 主要获奖情况}
\datedline{\textit{国家奖学金}}{2013}
\datedline{\textit{谷歌奖学金}}{2013}
\datedline{\textit{微软亚洲研究院实习生“明日之星”}}{2019}

%% Reference
%\newpage
%\bibliographystyle{IEEETran}
%\bibliography{mycite}
\end{document}
